\chapter{Tableaux croisés et fréquences conditionnelles}
\section*{Exercice 20.1 -- Compléter des effectifs}
Sur $200$ personnes, $120$ habitent le site A. Parmi elles, $90$ portent un casque. Sur le site B, $40$ portent un casque. Construire le tableau croisé, marges comprises.

\section*{Exercice 20.2 -- Choisir le dénominateur}
Avec ce tableau, calculer la proportion de porteurs de casque, la proportion de porteurs parmi les habitants de A et la proportion d'habitants de A parmi les porteurs. Interpréter chaque résultat.

\section*{Exercice 20.3 -- Nominal ou ordinal}
Classer les caractères suivants : canton de résidence ; satisfaction (faible, moyenne, forte) ; nombre d'enfants ; taille en centimètres. Pour lesquels une moyenne des modalités aurait-elle un sens ?

\section*{Exercice 20.4 -- Lire des données réelles}
Au 1\textsuperscript{er} janvier 2025, le canton d'Esch-sur-Alzette compte $97\,721$ femmes et $97\,370$ hommes. Le Luxembourg entier compte $339\,161$ femmes et $342\,812$ hommes (STATEC). Calculer la part des femmes dans ce canton et la part des résidentes du pays qui habitent ce canton.

\section*{Exercice 20.5 -- Filtrer avec ET, OU, NON}
Dans le tableau du casque, combien de personnes sont sur A ET portent un casque ? Sur A OU portent un casque ? Ni sur A ni porteuses de casque ? Expliquer pourquoi une addition directe échoue dans le deuxième cas.

\section*{Exercice 20.6 -- Une inversion trompeuse}
Une classe de $30$ élèves compte $18$ filles ; $12$ élèves font du théâtre, dont $9$ filles. « Les trois quarts des filles font du théâtre » est-il vrai ? Écrire l'énoncé réellement associé à $9/12$.
